\documentclass{article}
\usepackage[utf8]{inputenc}
\usepackage[english,ukrainian]{babel}
\usepackage{hyperref}
\usepackage{titlesec}
\usepackage{enumitem}

\titleformat{\section}{\large\bfseries}{}{0em}{}
\titleformat{\subsection}{\bfseries}{}{0em}{}

\begin{document}

\begin{center}
    {\LARGE\bfseries Сохацький Максим Еротейович}\\[0.4em]
    {\large namdak.tonpa@groupoid.space}\\[0.2em]
    \textbf{ORCID:} \href{https://orcid.org/0000-0001-7127-8796}{0000-0001-7127-8796}
\end{center}

\vspace{0.6em}

\begin{center}
    \small 26 листопада 1980 р., м. Кам'янець-Подільський, Хмельницька область
\end{center}

\vspace{0.8em}

\section{Наукові інтереси}
Теорія мов програмування, залежні типи (Dependent Type Systems), Homotopy Type Theory (HoTT),
вищі категорії і групоїди та їх семантика, категорна теорія та її застосування в програмуванні,
формальна верифікація програмного забезпечення, застосування теорії типів до задач інформаційної безпеки.

\section{Освіта}
\begin{itemize}[leftmargin=*, label={}]
    \item 1998--2005 — Національний технічний університет України «Київський політехнічний інститут» ім. Ігоря Сікорського, факультет прикладної математики.
          Ступінь: \textbf{магістр}.
    \item 2016--2022 — Аспірантура НТУУ «КПІ» ім. Ігоря Сікорського, факультет прикладної математики. \\
          Дисертаційну роботу завершено у 2022 році (захист відкладено).
    \item 1995--1997 — Приватний академічний ліцей «Антей» м. Кам'янець-Подільський (золота медаль).
\end{itemize}

\section{Професійний досвід}
\begin{itemize}[leftmargin=*]
    \item \textbf{Помічник директора}, Державне підприємство «Інформаційні судові системи», Судова влада України, Київ (квітень 2026 -- дотепер)
    \item \textbf{Директор Департаменту системного аналізу}, Державне підприємство «Електронне здоров'я», Міністерство охорони здоров'я України, Київ (червень 2024 -- березень 2026)
    \item \textbf{Enterprise Architect}, Державне підприємство «Інфотех», Міністерство внутрішніх справ України, Київ (вересень 2019 -- квітень 2023)
    \item \textbf{Старший науковий співробітник}, Національний технічний університет України «КПІ» ім. Ігоря Сікорського, Київ (вересень 2018 -- січень 2022)
    \item \textbf{Systems Architect}, АТ КБ «ПриватБанк», Дніпро (2014--2015)
%    \item \textbf{Software Developer}, ТОВ «ILS-Україна» (2003--2008)
%    \item \textbf{Software Developer}, ТОВ «Люксофт-Україна» (2010)
\end{itemize}

\section{Викладацька діяльність}
Автор курсів Інституту Формальної Математики «Групоїд Інфініті» (2023).

\subsection{Кафедра формальної філософії}
Теорія категорій, Формальні логіки, Категоріальна логіка, Лямбда-числення, Модальні логіки, Теорія типів Мартіна-Льофа, Формалізація штучних мов, Формалізація людських мов, Мультимодальна гомотопічна теорія типів, Теорія топосів, Теорія інфініті-категорій.

\subsection{Теорія мов програмування та суміжні дисципліни}
Теорія мов програмування, Історія мов програмування, Генератор компіляторів, Верифікація ПЗ, Верифікація об'єктного коду, Функціональне програмування, Лямбда-числення.

\subsection{Чиста та прикладна математика}
Гомотопічна теорія типів, Модельні категорії Квілена, Алгебраїчна топологія, Теорія гомотопій, Категоріальна логіка.

\section{Науково-дослідницька діяльність}
\begin{itemize}[leftmargin=*]
    \item \textbf{Незалежний дослідник та BDFL}, проект \textbf{Groupoid Space} та лабораторія \textbf{Groupoid Infinity} (2018 -- дотепер)
    \item Автор фундаментальної серії монографій \textbf{Groupoid Space} (8 томів)
    \item Керівник розробки академічної екосистеми \textbf{AXIO} — комплексу внутрішніх мов програмування та верифікаторів, побудованих на гомотопній теорії типів (HoTT), залежних типах та $\infty$-групоїдах
\end{itemize}

\textbf{Основні розроблені мови:}

\begin{itemize}[leftmargin=*]
    \item \textbf{alonzo, yves} — теоретичні інтерпретатори (декартово-замкнені та симетричні моноїдальні категорії)
    \item \textbf{henk, frank, christine, laurent, per, anders, dan, urs, fabien, jack} — математичні мови для механічного доведення теорем
    \item \textbf{tim, max, eijiro, manfred} — системні академічні мови
    \item \textbf{robin, leslie, andrea} — верифікатори програмного забезпечення та розподілених систем
\end{itemize}

Робота спрямована на уніфікацію синтетичної (HoTT, $\infty$-категорії) та класичної математики в єдиному механічно верифікованому фреймворку. Всі проєкти open-source, доступні на \href{https://github.com/groupoid}{github.com/groupoid}.

\section{Книги}
Автор фундаментальної серії \textbf{Groupoid Space}, що складається з 8 томів:

\begin{itemize}[leftmargin=*]
    \item \textbf{Volume I} — Foundations
    \item \textbf{Volume II} — Theoretical Interpreters
    \item \textbf{Volume III} — Fibrational Provers
    \item \textbf{Volume IV} — Advanced Studies
    \item \textbf{Volume V} — Verification
    \item \textbf{Volume VI} — Formal Philosophy
    \item \textbf{Volume VII} — Introduction to Mind
    \item \textbf{Volume VIII} — Philosophy of Mathematics
\end{itemize}

Всі томи доступні у відкритому доступі на \href{https://groupoid.space}{https://groupoid.space}.

\section{Статті та публікації}
Основні статті входять до складу серії \textbf{Groupoid Space} (том IV та інші).

\subsection{Основи}
Type Theory, Inductive Types, Homotopy Type Theory, Higher Inductive Types, Modalities.

\subsection{Мови}
Alonzo Church, Yves Lafont, Felix Bloch, Joe Armstrong, Robin Milner, Henk Barendregt, Per Martin-Löf, Frank Pfenning, Laurent Schwartz, Anders Mörtberg, Urs Schreiber.

\subsection{Математики}
Algebra vs Geometry, Category Theory, Topos Theory, Categories with Families, Categories with Representable Maps, Abelian Categories, Comprehension Categories, Cosmic Cube, Homotopical Dynamics, Fibered Categories, Chevalley Descent, Local Cartesian Closed Categories, Symmetric Monoidal Categories, Cohesive Topoi, Quillen Model Structure, Grothendieck Schemes, Simplicial Homotopy Theory, Cohomology and Spectra, Grothendieck Yoga.

\section{Людські мови}
\begin{itemize}[leftmargin=*]
    \item Українська — рідна
    \item Англійська — C1 (науковий рівень)
    \item Тибетська — дхармічна мова, зі словником
\end{itemize}

\section{Додаткова інформація}
\begin{itemize}[leftmargin=*]
    \item Проживаю: 04071, м. Київ, вул. Костянтинівська, 20, кв. 31
    \item ФОП з 2010 року
    \item Директор релігійної організації «Лонгчен Нінгтік» (ЄДРПОУ 38778275)
    \item Бенефеціар ТОВ «Криптографічні Телесистеми» (ЄДРПОУ 46226772)
\end{itemize}

\vspace{1.2em}
\begin{center}
    \small Київ, травень 2026 року
\end{center}

\end{document}
